\section{Discussion}
\label{sec:discussion}

\subsection{What the polynomial activation actually contributes}

Taken at face value, Table~\ref{tab:activation} says that $z + \lambda z^{3}$
is a \emph{weaker} activation than the modern defaults: at depth ten, width
$128$ and equal parameter count the bare cubic lands almost four accuracy
points below ReLU, GELU and SiLU, and the $\ell_2$ projection closes a little
over half of that distance without closing it. No amount of polynomial degree
or coefficient tuning removes the remainder. That reading is worth stating
plainly because activation-function
folklore often promises more. On a task whose decision boundaries are
essentially smooth, a bounded-degree polynomial purchases no accuracy advantage
over a piecewise-linear or a smooth saturating nonlinearity, and in our
measurements it pays a penalty instead.

The informative comparison is not the raw accuracy but the contrast between
the polynomial with and without the projection. Adding the $\ell_2$ projection
changes the result far more than changing the polynomial's degree or its
coefficient does --- it is worth $+\polyGap$ points, roughly half the deficit to
ReLU. This is the paper's central empirical claim, and it re-frames the
original construction: what the projective polynomial offers is not a better
nonlinearity, but a \emph{parameter-free, batch-independent mechanism for
maintaining activation scale across depth}.

\subsection{Why the projection substitutes for BatchNorm and LayerNorm}

BatchNorm normalizes using statistics estimated across the
batch~\citep{ioffe2015batch};
LayerNorm normalizes using statistics estimated across features within each
example~\citep{ba2016layer}. Both therefore add a learned affine
transformation and a pair of
statistics buffers per layer, and both change the effective function computed
by the network in a data-dependent way. The $\ell_2$ projection does neither:
it rescales each activation vector to unit norm, which is exactly a
LayerNorm-like operation with the mean subtraction and the affine map removed.

Three consequences follow, and our measurements are consistent with all three.
First, the projection is parameter-free: it adds no learned parameters to any
layer. Second, it carries no cross-example coupling, so it does not require a
batch dimension to be meaningful. Third, and this is where the depth sweep
matters, it preserves the ability to train at depths where the unnormalized
nonlinearity fails, and it does so at accuracy comparable to the normalized
baselines without their parameter overhead.

\subsection{Gated units: a parameter-budget question, not an accuracy question}

The gated results deserve a careful reading, because it would be easy to
conclude from them either that gating helps or that it hurts, and both
conclusions would be wrong.

When parameter counts are matched, the advantage of the gated variants largely
disappears. That is the expected outcome, and it is the reason we compare at
equal parameter count rather than equal width: a gated unit at width $w$ has
two projections where a pointwise activation has one, so comparing at equal
width silently gives the gate roughly three times the parameters. Once that
is controlled, the gate is not obviously the source of any accuracy gain.

More strikingly, a gated unit does not train at depth ten on MNIST without
normalization at all. We verified that this is a real numerical failure rather
than a bug: the activation norm decays by many orders of magnitude across the
stack, the output gradient is exactly zero, and the network remains at chance.
Adding normalization of any kind --- $\ell_2$, batch, layer, or RMS ---
restores training immediately. This is a clean demonstration of a general
point that gating is not a free-standing architectural improvement: a
multiplicative interaction between two unnormalized signals produces a scale
that compounds multiplicatively with depth, and only an explicit
renormalization keeps that compounding from driving the signal to zero or to
overflow. In large language models, gated units appear inside residual blocks
that already carry normalization, so the failure mode is masked in practice;
our experiment isolates it.

\subsection{Limitations}

Several limitations should be kept in view when reading the numbers.

\emph{Task family.} The main benchmark is MNIST, whose classes are separated
by smooth, low-dimensional decision boundaries. It is a well-controlled
setting for isolating optimization effects, but it is not a setting in which a
polynomial activation can show a representational advantage --- indeed, our
synthetic benchmarks show that even on curved targets (rings, spirals, the
checkerboard) the polynomial does not pull ahead of the modern elementwise
activations. Those synthetic results should be read as indicative rather than
conclusive, since each synthetic dataset uses a single train/test split.

\emph{Architecture.} We study fully-connected stacks. Our conclusions about
the projection as a substitute for normalization are statements about
scale control in a feed-forward stack; they do not automatically transfer to
convolutional or attention architectures, where normalization interacts with
residual connections and with the attention softmax in ways this study does
not probe.

\emph{Optimizer and schedule.} All runs use Adam with a fixed learning rate
and a fixed training budget. Because normalization changes the geometry of the
loss surface, it is plausible that the normalization comparison would shift
somewhat under a tuned per-configuration learning rate or a longer schedule.
We deliberately hold the optimizer fixed so that the comparison reflects
architecture rather than tuning effort; a per-configuration sweep is a natural
extension.

\emph{Statistics.} The core comparisons use a single seed. The multi-seed
control over the principal configurations shows the seed-to-seed spread to be
small relative to the differences we report as meaningful, and the conclusions
rest on differences of one to three accuracy points together with a
qualitative training failure, not on differences at the level of the noise.
